% source: manti.01.ma, 1ère SM, série fonctions (Drive 0B6pfSJcEVng0WHNzSmtjS0xudk0).
% Ex.1 (1,2,4,5,6) et Ex.2 (1,3) = série française fonctions_s1 -> non repris
\ex{Fonctions}{WHNz}{Ex. 1}{6}
\begin{enumerate}
\item Soit $f$ la fonction définie par $f(x)=x\sqrt{1+x^2}-x^2$. Montrer que $f$ est majorée par $\dfrac12$.
\item Soit $f$ la fonction définie par $f(x)=\dfrac{x^2+4x+1}{x^2+1}$.
\begin{enumerate}
\item Déterminer l'ensemble de définition $D_f$.
\item Montrer que $f$ admet un maximum absolu en $1$ et un minimum absolu en $-1$.
\end{enumerate}
\item Soit $f$ la fonction définie sur $\mathbb{R}$ par $f(x)=\dfrac{|x|}{x^2+2}$.
\begin{enumerate}
\item Montrer que $f(x)<1$ pour tout $x$ de $\mathbb{R}^+$.
\item En déduire que $f$ est bornée.
\item Montrer que $\max f(x)=\dfrac{\sqrt2}{4}$.
\end{enumerate}
\end{enumerate}
\fin
% Ex. 3 = série française fonctions_s2 (non repris)
\ex{Fonctions}{WHNz}{Ex. 4}{7}
On pose $f(x)=\dfrac{x^2}{x^2+2}$.
\begin{enumerate}
\item Étudier la parité de la fonction $f$.
\item Montrer que $f$ est majorée par $1$.
\item Montrer que $\min f(x)=f(0)$.
\item Étudier les variations de la fonction $f$.
\item Montrer que $f(\mathbb{R})=[0,1[$.
\end{enumerate}
\fin
\ex{Fonctions}{WHNz}{Ex. 5}{8}
Soit $f$ la fonction définie par $f(x)=\sqrt x+\sqrt{10-x}$.
\begin{enumerate}
\item Déterminer l'ensemble de définition $D_f$.
\item Montrer que pour tout $x$ de $D_f$, $10-x$ appartient à $D_f$, puis calculer $f(10-x)$.
\item Montrer que $f$ est croissante sur $[0,5]$ et en déduire la monotonie de $f$ sur $[5,10]$.
\item Comparer les nombres $\sqrt3+\sqrt7$ et $\sqrt2+\sqrt8$.
\end{enumerate}
\fin
